\documentclass[10pt,a4paper]{amsart}
\usepackage[margin=19mm]{geometry}
\usepackage{amsmath,amssymb,mathtools}
\usepackage[scr=rsfs,cal=euler]{mathalfa}
\usepackage{xfrac}
\usepackage[unicode,hidelinks,bookmarks=false]{hyperref}
\newcommand{\ie}{i.e.\ }
\newcommand{\cf}{cf.\ }
\newcommand{\resp}{respectively\ }
\newcommand{\Subsection}{\subsection}
\providecommand{\nobreakdash}{\nobreak\hbox{-}}
\setlength{\emergencystretch}{2em}
\multlinegap0pt
\usepackage{amssymb}
\usepackage{mathtools}
\usepackage{dsfont}
\usepackage{algorithm}\usepackage{algpseudocode}
\usepackage{tikz}
\usepackage[shortlabels]{enumitem}
\usepackage{xcolor}
\newtheorem{prop}{Proposition}
\newcommand{\bbm}{\begin{bmatrix}}
\newcommand{\ebm}{\end{bmatrix}}
\newcommand{\re}[1]{\mathbb{R}^{#1}}%
\title{On maximal positive invariant set computation for rank-deficient linear systems}
\author{Bogdan Gheorghe, Daniel Ioan, Cristian Flutur, Ionela Prodan,, Florin Stoican}
\date{Source: arXiv:2603.17686v1 (CC BY 4.0)}
\begin{document}
\maketitle

\noindent\emph{Source and licence.} On maximal positive invariant set computation for rank-deficient linear systems. Version arXiv:2603.17686v1 (CC BY 4.0), distributed by arXiv under the Creative Commons Attribution 4.0 International licence, http://creativecommons.org/licenses/by/4.0/, as recorded in the arXiv licence metadata for this exact version and shown on the abstract page https://arxiv.org/abs/2603.17686v1. Full licence notice: \texttt{licenses/arxiv-2603.17686-CC-BY-4.0.txt}.\par\smallskip

\noindent\textit{Editorial scope.} Counted unit: the article's prop prop:cz (source0.tex lines 614--645). The article's complete original proof of it is delivered with it (source0.tex lines 646--658, one \texttt{proof} environment of 13 lines). No mathematics is written, repaired or re-derived: the statement and the proof are the article's own lines, in the article's own notation, and no retained line is substituted or re-flowed.  Retained uncounted as the source-local context the proof needs: lines 330--333; lines 360--363; lines 365--369; lines 508--514; lines 516--519; lines 523--526; lines 557--564; lines 568--595.  Nothing was omitted from any retained span, so the package carries no omission record.  The retained lines cite 1 works, whose entries are transcribed verbatim from the article's own bibliography source in the e-print and delivered with short editorial labels recorded in the item's reference mapping.  No retained line contains a figure environment, an image-inclusion command or drawing code, so no image is delivered and the package declares no asset dependency.  Editorial formatting only, in addition: an AMS \texttt{amsart} preamble that restates the article's own declarations and macros used by the retained lines, namely the environments \texttt{prop} on separate counters, together with the standard packages those lines need.  The retained environments print in this document as Proposition 1 in order of appearance, whereas the article prints the same items with the numbering of its own full document.  The complete native source of the article as distributed in the arXiv e-print for this exact version is bundled unchanged as \texttt{sources/196639/source0.tex}.
\begin{equation}
\label{eq:con_poly_zono}
    \mkern-8mu\mathcal{CZ} = \bigl\langle c, G, F, \theta\bigr\rangle = \big \{x = c + G\lambda,\: F\lambda=\theta,\: \|\lambda\|_\infty \leq 1\bigr\}.
\end{equation}

\begin{multline}
\label{eq:x_constr_zon}
\overline{\mathcal X} = \left\langle c, G, F, \theta\right\rangle =\bigl\langle c, \bbm G& 0\ebm,\\ \bbm KG& -G_U\ebm, c_U-Kc\bigr\rangle.
\end{multline}

\begin{multline}
\label{eq:mpi_constr_zon}
\Omega_{k+1} =\Bigl\langle A_\circ^{-1}c_k, \bbm A_\circ^{-1}G_k& 0\ebm,\\
    \bbm F_k&0\\ 0& F\\ A_\circ^{-1}G_k& -G\ebm,  \bbm\theta_k\\\theta\\c-A_\circ^{-1}c_k\ebm\Bigr\rangle.
\end{multline}

\begin{equation}
\label{eq:schur}
    A_\circ=A+BK = U\left[\begin{tabular}{c|c}
         $S_{11}$& $S_{12}$ \\\hline
         0& $S_{22}$
    \end{tabular}\right]U^\top, 
\end{equation}

\begin{equation}
\label{eq:nilpotent}
    \bigl(S_{22}\bigr)^{p+1} = 0, \qquad \text{eig}(S_{22}) = 0.
\end{equation}

\begin{equation}
    \label{eq:T}
    T=\sum_{i=0}^{p} \bigl(S_{11}\bigr)^{-i}S_{12}\bigl(S_{22}\bigr)^{i},
\end{equation}

\begin{align}
    \label{eq:y_dynamics}
    y_{k+1}&=\left[\begin{tabular}{c|c}
         $S_{11}$& $S_{12}$ \\\hline
         0& $S_{22}$
    \end{tabular}\right]y_k, \\
    \label{eq:y_set}\mathcal Y &= \{y:\: FUy\leq \theta\}.
\end{align}

Decomposing $y=\bbm y_1^\top & y_2^\top \ebm^\top$ we use \eqref{eq:y_dynamics} to write the evolution of each component of the state
\begin{subequations}
\label{eq:dynamics}
    \begin{align}
        \nonumber y_{1, k+1} &= S_{11}\,y_{1, k} + S_{12}y_{2, k}\\
        &=\bigl(S_{11}\bigr)^{k+1}y_{1,0}+\sum\limits_{i=0}^k \bigl(S_{11}\bigr)^{k-i}S_{12}\bigl(S_{22}\bigr)^{i}y_{2,0},\\
        \nonumber y_{2, k+1} &= S_{22}y_{2, k}\\
        &=\bigl(S_{22}\bigr)^{k+1}y_{2,0}.
    \end{align}
\end{subequations}
% For further use we make the shorthand notation $T_{k,i}=\bigl(S_{11}\bigr)^{k-i}S_{12}\bigl(S_{22}\bigr)^{i}$. 
Exploiting the nilpotence condition \eqref{eq:nilpotent} we rewrite \eqref{eq:dynamics} into
\begin{subequations}
\label{eq:dynamics_nilpotent}
    \begin{align}
        \nonumber y_{1, k+1} =&  \bigl(S_{11}\bigr)^{k+1}y_{1,0}\\
        &+          \label{eq:dynamics_nilpotent_a}\bigl(S_{11}\bigr)^{k}\sum\limits_{i=0}^{\min(k,\, p)} \bigl(S_{11}\bigr)^{-i}S_{12}\bigl(S_{22}\bigr)^{i}y_{2,0},\\
        \label{eq:dynamics_nilpotent_b}y_{2, k+1} =& \begin{cases}\bigl(S_{22}\bigr)^{k+1}y_{2,0}, & k<p,\\
        0, & k\geq p.
        \end{cases}
    \end{align}
\end{subequations}
% Clearly, for any $k\geq p$, component $y_{2,k+1}=0$ as per \eqref{eq:dynamics_nilpotent_b}, and only $y_{1,k+1}$ updates as in \eqref{eq:dynamics_nilpotent_a}.
For convenience, let us denote $z_k=y_{1,k+p}$ and take $T$ as in \eqref{eq:T}. Consequently, for any $k\geq p$, \eqref{eq:dynamics_nilpotent_a} is equivalently written as (we shift from $k$ to $k-p$)
\begin{equation}
\label{eq:z}
    z_{k+1} = S_{11} z_k, \quad z_0=\bigl(S_{11}\bigr)^py_{1,0}+\bigl(S_{11}\bigr)^{p-1}Ty_{2,0}.
\end{equation}

\begin{prop}
\label{prop:cz}
    Consider the notation from \eqref{eq:schur}--\eqref{eq:T} and the constraint set given in its constrained zonotopic form $\overline{\mathcal X}=\langle c,G,F,\theta\rangle$, defined as in \eqref{eq:x_constr_zon}. Then the associated MPI set is given by 
    \begin{multline}
    \label{eq:mpi-schur-cz}
        \mkern-36mu\Phi_x = \mkern-2mu\biggl\langle \mkern-4mu c^{[1]},\mkern-2mu \bbm G^{[1]} & 0\ebm\mkern-2mu,\mkern-4mu \bbm F^{[1]} & 0\\ 0 & F^{[2]}\\ \bbm \bigl(S_{11}\bigr)^p &\bigl(S_{11}\bigr)^{p-1}T\ebm U^\top G^{[1]} & -G^{[2]}\mkern-4mu\ebm\mkern-4mu,\\
        \bbm \theta^{[1]}\\ \theta^{[2]}\\ c^{[2]}-\bbm \bigl(S_{11}\bigr)^p &\bigl(S_{11}\bigr)^{p-1}T\ebm c^{[1]}\ebm\biggr\rangle
    \end{multline}
    where the ``$[1]$'' parameters describe the constraints imposed until $k<p$:
    \begin{equation}
    \label{eq:omega_x_1_cz}
        \Phi_x^{[1]}=\bigcap\limits_{k=0}^{p-1} \overline{\mathcal X} = \bigl\langle c^{[1]}, G^{[1]}, F^{[1]}, \theta^{[1]}\bigr\rangle
    \end{equation}
and the ``$[2]$'' parameters describe the constraints applied to $k\geq p$ and, under dynamics \eqref{eq:z} and set recurrence \eqref{eq:mpi_constr_zon} which stops after $\bar k$ iterations \cite{10552328}:
    \begin{multline}
        \label{eq:omega_x_2_cz}
        \Phi_x^{[2]}=\bigcap\limits_{k=p}^{\bar k +p} {\mathcal X}\\
        =\{x:\: \bbm \bigl(S_{11}\bigr)^p &\bigl(S_{11}\bigr)^{p-1}T\ebm U^\top x\in \Phi_z\}, \theta^{[2]}\bigr\rangle.
    \end{multline}
    where
    \begin{equation}
        \label{eq:mpi_z_cz}
        \Phi_z =\bigcap\limits_{k=0}^{\bar k} {\mathcal Z}=\langle c^{[2]}, G^{[2]}, F^{[2]}, \theta^{[2]}\rangle,
    \end{equation}
    is the MPI set associated with the dynamics $z^+=S_{11}z$ and constraint set 
    \begin{equation}
    \label{eq:z-cz}
        \mathcal Z = \biggl\langle U_1^\top c, U_1^\top G, \bbm F\\ U_2^\top G\ebm, \bbm \theta\\ -U_2^\top c\ebm\biggr\rangle.
    \end{equation}
    % where, in \eqref{eq:schur}, $U= \bbm U_1 & U_2\ebm$, $U_1 \in \re{n\times d_1}$, $U_2 \in \re{n\times d_2}$. 
    All $(\cdot)^{[1]},(\cdot)^{[2]}$ parameters are obtained recursively as in \eqref{eq:mpi_constr_zon}.
\end{prop}

\begin{proof}
Using recurrence \eqref{eq:mpi_constr_zon} until step $k=p-1$ for closed-loop matrix \eqref{eq:schur} and constraint set \eqref{eq:x_constr_zon} gives $\Phi_x^{[1]}$ as in \eqref{eq:omega_x_1_cz}. Next, for $k\geq p$, we consider the same change of coordinates as before, $y=U^\top x$. Using definition \eqref{eq:con_poly_zono}, we have that $x_{k+1}\in \overline{\mathcal X}$ can be written as
\begin{subequations}
\label{eq:y_cz}
\begin{align}
\label{eq:y_cz_1}
    y_{1, k+1}&=U_1^\top c + U_1^\top G\lambda_{k+1},\quad F_1\lambda_{k+1} = \theta_1,\\
\label{eq:y_cz_2}
y_{2, k+1}&=U_2^\top c + U_2^\top G\lambda_{k+1},\quad F_2\lambda_{k+1} = \theta_2.
\end{align}    
\end{subequations}
Due to \eqref{eq:nilpotent} and as per \eqref{eq:dynamics_nilpotent_b} we have that $y_{2, k+1}=0$. This acts as an additional equality constraint on $\lambda_{k+1}$ in \eqref{eq:y_cz_2}. Introducing it into \eqref{eq:y_cz_1}, leads to inclusion $y_{1,k+1}\in \mathcal Z$, with $\mathcal Z$ defined as in \eqref{eq:z-cz}. With notation \eqref{eq:z} and the set \eqref{eq:z-cz}, set recurrence \eqref{eq:mpi_constr_zon} leads to \eqref{eq:mpi_z_cz}. Using the second part of \eqref{eq:z} and \eqref{eq:omega_x_1_cz} and \eqref{eq:omega_x_2_cz} leads to \eqref{eq:mpi-schur-cz}, thus concluding the proof.
\end{proof}

\begin{thebibliography}{99}
\bibitem[105523]{10552328} B.~Gheorghe, D.-M. Ioan, F.~Stoican, and I.~Prodan, ``Computing the maximal positive invariant set for the constrained zonotopic case,'' \emph{IEEE Control Systems Letters}, vol.~8, 2024.
\end{thebibliography}
\end{document}
